\documentclass[exercice]{thedocument}%
\startdatakeys
\source{}
\BO{2026}
\cycle{4}
\competencesDuSocle{CA,CO}
\objectifsApprentissage{A2ab,A1ak,A1bg}
\automatismes{B1eb,A1cf}
\enddatakeys
\begin{document}%
Trois élèves ont calculé l'expression numérique suivante :
$$A=3+2\times5-4$$
Léa a trouvé 5, Christine 2 et Sam 9.\vskip5pt
\begin{enumerate}
    \item Lequel des trois a obtenu la réponse correcte ?
    \item Expliquer l'erreur commise par chacun des deux autres.
\end{enumerate}
\end{document}
